% Options for packages loaded elsewhere
\PassOptionsToPackage{unicode}{hyperref}
\PassOptionsToPackage{hyphens}{url}
%
\documentclass[
  11pt,
]{article}
\usepackage{amsmath,amssymb}
\usepackage{iftex}
\ifPDFTeX
  \usepackage[T1]{fontenc}
  \usepackage[utf8]{inputenc}
  \usepackage{textcomp} % provide euro and other symbols
\else % if luatex or xetex
  \usepackage{unicode-math} % this also loads fontspec
  \defaultfontfeatures{Scale=MatchLowercase}
  \defaultfontfeatures[\rmfamily]{Ligatures=TeX,Scale=1}
\fi
\usepackage{lmodern}
\ifPDFTeX\else
  % xetex/luatex font selection
\fi
% Use upquote if available, for straight quotes in verbatim environments
\IfFileExists{upquote.sty}{\usepackage{upquote}}{}
\IfFileExists{microtype.sty}{% use microtype if available
  \usepackage[]{microtype}
  \UseMicrotypeSet[protrusion]{basicmath} % disable protrusion for tt fonts
}{}
\makeatletter
\@ifundefined{KOMAClassName}{% if non-KOMA class
  \IfFileExists{parskip.sty}{%
    \usepackage{parskip}
  }{% else
    \setlength{\parindent}{0pt}
    \setlength{\parskip}{6pt plus 2pt minus 1pt}}
}{% if KOMA class
  \KOMAoptions{parskip=half}}
\makeatother
\usepackage{xcolor}
\usepackage[margin=21mm]{geometry}
\usepackage{longtable,booktabs,array}
\usepackage{calc} % for calculating minipage widths
% Correct order of tables after \paragraph or \subparagraph
\usepackage{etoolbox}
\makeatletter
\patchcmd\longtable{\par}{\if@noskipsec\mbox{}\fi\par}{}{}
\makeatother
% Allow footnotes in longtable head/foot
\IfFileExists{footnotehyper.sty}{\usepackage{footnotehyper}}{\usepackage{footnote}}
\makesavenoteenv{longtable}
\setlength{\emergencystretch}{3em} % prevent overfull lines
\providecommand{\tightlist}{%
  \setlength{\itemsep}{0pt}\setlength{\parskip}{0pt}}
\setcounter{secnumdepth}{-\maxdimen} % remove section numbering
\ifLuaTeX
\usepackage[bidi=basic]{babel}
\else
\usepackage[bidi=default]{babel}
\fi
\babelprovide[main,import]{french}
% get rid of language-specific shorthands (see #6817):
\let\LanguageShortHands\languageshorthands
\def\languageshorthands#1{}
\ifLuaTeX
  \usepackage{selnolig}  % disable illegal ligatures
\fi
\IfFileExists{bookmark.sty}{\usepackage{bookmark}}{\usepackage{hyperref}}
\IfFileExists{xurl.sty}{\usepackage{xurl}}{} % add URL line breaks if available
\urlstyle{same}
\hypersetup{
  pdftitle={Signe et couverture : le coup double précisément situé},
  pdfauthor={Vincent P. P. NOËL --- dossier de recherche, rédaction assistée},
  pdflang={fr},
  hidelinks,
  pdfcreator={LaTeX via pandoc}}

\title{Signe et couverture : le coup double précisément situé}
\usepackage{etoolbox}
\makeatletter
\providecommand{\subtitle}[1]{% add subtitle to \maketitle
  \apptocmd{\@title}{\par {\large #1 \par}}{}{}
}
\makeatother
\subtitle{Complément au document 76 --- Programme Riemann V6.55}
\author{Vincent P. P. NOËL --- dossier de recherche, rédaction assistée}
\date{15 octobre 2026}

\begin{document}
\maketitle

\hypertarget{ce-qui-change-dans-la-direction-de-recherche}{%
\section{1. Ce qui change dans la direction de
recherche}\label{ce-qui-change-dans-la-direction-de-recherche}}

\textbf{Le transfert de couverture existe déjà dans le corpus. Le
mécanisme arithmétique qui fournirait le signe global reste ouvert.} Le
document 67, pont B3, construit un cœur dense en norme \(H^1\) à support
fixé et justifie le transfert d'une positivité établie sur toute la
famille. Le document 76 donne des grammaires couvrantes : paquets
traduits et deux filtres. Il serait inexact de remettre ces résultats à
zéro.

La proposition « trouver \(Q(f)=\|Tf\|^2\) » reçoit ici un opérateur
candidat concret. Le résultat de cette passe est un contrat de travail
vérifiable : établir une identité de Gram sur un cœur dense, puis
employer une continuité obtenue indépendamment de cette identité.
L'identité de Gram n'est pas prouvée ici.

Ce complément conserve les numéros des documents, les identifiants des
idées et les outils de pilotage. Il prolonge les fiches IDE-044
(domaines), IDE-047 (Schur), IDE-050 (couverture), IDE-051 (signe) et
IDE-052 (Weil). Les preuves ci-dessous sont des résultats auxiliaires,
sans revendication de nouveauté et sans validation extérieure.

\hypertarget{lopuxe9rateur-de-la-littuxe9rature-et-sa-garde-de-continuituxe9}{%
\section{2. L'opérateur de la littérature et sa garde de
continuité}\label{lopuxe9rateur-de-la-littuxe9rature-et-sa-garde-de-continuituxe9}}

{[}SuzH23{]} construit sans RH les fonctions
\(\mathfrak S_t\in L^2(\mathbb R)\) et l'opérateur \[
T f=\pi^{-1/2}\widehat{\mathcal P_{Df}},\qquad Df=if'.
\] Les propositions 1.2--1.3 définissent ce candidat ; le théorème 1.4
rend l'identité \(\|Tf\|_2^2=\langle f,f\rangle_W\) sur tous les tests
lisses compacts équivalente à RH. Le §3.2 établit aussi une borne
uniforme de \(\|\mathfrak S_t\|_2\) sur chaque compact de \(t\). Nous
utilisons ces résultats publiés comme entrées, sans redémontrer leur
analyse complexe.

\textbf{Convention.} Nous notons \(q(f,g)\) la forme de Weil de cette
référence, rendue antilinéaire en première variable. Sa valeur diagonale
reste celle de la source. Le \(Q_W\) arithmétique du document 76 emploie
sa propre convention Fourier et des moments polaires nuls. Avant tout
calcul comparatif, il faut dériver leur concordance exacte à partir de
la formule explicite, facteur de normalisation compris. Ce raccord n'est
pas certifié numériquement dans cette passe.

\textbf{Lemme A --- Continuité indépendante de l'identité.} Pour
\(S>0\), poser \[
M_S=\sup_{|t|\le S}\|\mathfrak S_t\|_2<\infty,
\qquad C_S=\sqrt{2S/\pi}\,M_S.
\] Alors, pour \(f\in C_c^\infty(-S,S)\), \[
\boxed{\|Tf\|_2\le C_S\|f'\|_2.}
\] \textbf{Preuve.} La formule intégrale de \(\widehat{\mathcal P}\)
donne \[
\|Tf\|_2\le\pi^{-1/2}\int_{-S}^{S}
 \|\mathfrak S_t^\sharp\|_2|f'(t)|\,dt
\le\pi^{-1/2}M_S\|f'\|_1
\le C_S\|f'\|_2.
\] Sur l'axe réel, la conjugaison \(\sharp\) préserve la norme.
Minkowski puis Cauchy--Schwarz donnent les deux inégalités. L'opérateur
s'étend donc par continuité à \(H_0^1(-S,S)\). Les extensions pour
supports emboîtés coïncident sur leur domaine commun, par densité et
unicité. Aucune constante uniforme en \(S\) n'est nécessaire. Aucune
positivité de Weil n'a été utilisée.

\textbf{Portée.} Ce lemme raccorde le candidat publié à la topologie de
couverture du corpus. La constante est finie ; aucune valeur numérique
certifiée de \(M_S\) n'est fournie ici.

\hypertarget{le-duxe9faut-qui-porte-lobligation-globale}{%
\section{3. Le défaut qui porte l'obligation
globale}\label{le-duxe9faut-qui-porte-lobligation-globale}}

Dans un espace \(E_S\) où \(q\) est continu en norme \(H^1\), définir \[
\mathfrak D(f,g)=q(f,g)-\langle Tf,Tg\rangle_2,
\qquad d(f)=\mathfrak D(f,f).
\] On peut choisir le sous-espace à moments polaires nuls du document
67, après raccord des normalisations. Pour le critère publié sur tous
les tests compacts, choisir le domaine de forme correspondant et sa
continuité. Il faut déclarer lequel est employé ; ces choix ne se
substituent pas silencieusement.

\textbf{Lemme B --- Identité sur un cœur, signe sur le domaine.} Soient
\(f_j\) tels que l'union de
\(V_n=\operatorname{span}_{\mathbb C}(f_1,\ldots,f_n)\) soit dense dans
\(E_S\). Si \[
\mathfrak D(f_i,f_j)=0\quad\text{pour tous }i,j,
\] alors \(q(f,f)=\|Tf\|_2^2\ge0\) pour tout \(f\in E_S\).

\textbf{Preuve.} La sesquilinéarité donne la nullité sur chaque
\(V_n\times V_n\). Pour \(f_n\to f\) dans \(H^1\), la continuité de
\(q\) donne \(q(f_n,f_n)\to q(f,f)\) et le lemme A donne \(Tf_n\to Tf\)
dans \(L^2\). Passer à la limite dans l'identité prouve le résultat.
Traiter chaque support entier \(S\ge1\) couvre tout test compact. Cette
implication n'exige aucune minoration strictement positive uniforme.

\textbf{Variante moins exigeante.} Une factorisation exacte est
suffisante, mais elle n'est pas nécessaire pour établir le signe. Si,
pour tous \(n\), \[
q(v,v)\ge-\varepsilon_{S,n}\|v\|_{H^1}^2
\quad(v\in V_n),\qquad\varepsilon_{S,n}\longrightarrow0,
\] le signe suit également. En effet, un \(v\in V_m\) appartient à tous
les \(V_n\) ultérieurs ; sa norme reste fixe, puis la densité achève la
preuve. Ce transfert était déjà présent dans le document 67. Il reste à
produire l'estimation arithmétique pour tous les rangs.

\hypertarget{pourquoi-il-faut-tester-les-muxe9langes}{%
\section{4. Pourquoi il faut tester les
mélanges}\label{pourquoi-il-faut-tester-les-muxe9langes}}

Avec notre convention antilinéaire en première variable, \[
\operatorname{Re}\mathfrak D(f,g)
=\frac{d(f+g)-d(f-g)}4,
\] \[
\operatorname{Im}\mathfrak D(f,g)
=-\frac{d(f+ig)-d(f-ig)}4.
\] Ces identités suivent en développant les quatre formes quadratiques.
Elles déterminent exactement les couplages et leurs phases.

\textbf{Contre-exemple réel.} Pour \(T=I\) sur \(\mathbb C^2\) et \[
Q=\begin{pmatrix}1&2\\2&1\end{pmatrix},
\] les deux vecteurs de base ont la bonne énergie, mais \(q(1,-1)=-2\).
Une vérification brique par brique échoue.

\textbf{Contre-exemple complexe.} Pour \[
Q=\begin{pmatrix}1&2i\\-2i&1\end{pmatrix},
\] \(q(x)=\|x\|^2\) pour tout vecteur réel \(x\), tandis que
\(q(1,i)=-2\). Même tous les tests réels ne valident pas une forme
hermitienne arbitraire. Une symétrie particulière peut réduire les tests
seulement après démonstration de cette symétrie.

Ces deux matrices sont des contre-modèles abstraits de protocoles
insuffisants ; elles ne sont pas des matrices de Weil. Pour l'opérateur
exact de Suzuki, un défaut non nul démontré rigoureusement aurait une
portée sur RH. Une différence flottante sans contrôle des intégrales,
singularités apparentes et queues ne suffit pas.

\hypertarget{la-brique-dassemblage-schur-sans-perdre-les-noyaux}{%
\section{5. La brique d'assemblage : Schur sans perdre les
noyaux}\label{la-brique-dassemblage-schur-sans-perdre-les-noyaux}}

Soit \(M\succeq0\) une matrice hermitienne, \(b\) une colonne et \(a\)
un réel. Poser \[
\widetilde M=\begin{pmatrix}M&b\\b^*&a\end{pmatrix}.
\] \textbf{Lemme C.} Avec \(M^\dagger\) la pseudo-inverse de
Moore--Penrose, \[
\widetilde M\succeq0
\iff b\in\operatorname{Ran}M\ \text{et}\
\delta=a-b^*M^\dagger b\ge0.
\] \textbf{Preuve.} Si \(z\in\ker M\), la restriction à \((tz,1)\)
impose \(z^*b=0\), sinon le terme linéaire prend des valeurs négatives
arbitraires. Ainsi \(b\perp\ker M\), donc \(b\in\operatorname{Ran}M\).
Sous cette condition, compléter le carré sur l'image donne \[
x^*Mx+2\operatorname{Re}(x^*by)+a|y|^2
=(x+M^\dagger by)^*M(x+M^\dagger by)+\delta|y|^2.
\] La nécessité et la suffisance suivent.

Si \(M=V^*V\), on garde les anciennes colonnes et on ajoute \[
w=VM^\dagger b+\sqrt\delta\,e,
\] où \(e\) est un vecteur unitaire orthogonal à toutes les anciennes
colonnes. On obtient \(V^*w=b\) et \(\|w\|^2=a\) : la réalisation Gram
est compatible avec l'étape précédente.

\textbf{Usage et limite.} Une règle arithmétique qui impose
l'appartenance à l'image et \(\delta\ge0\) à tous les rangs ferait le
coup double sur une famille couvrante. Le lemme ne fournit pas cette
règle. Factoriser des matrices dont on connaît déjà le signe ne
constitue pas une preuve supplémentaire de ce signe. La garde d'image
reste indispensable lorsqu'un pivot s'annule.

\hypertarget{deux-fausses-couvertures-uxe9cartuxe9es}{%
\section{6. Deux fausses couvertures
écartées}\label{deux-fausses-couvertures-uxe9cartuxe9es}}

\textbf{Mauvaise norme.} Sur \(C_c^\infty(-1,1)\), poser
\(r(f)=\|f\|_2^2-|f(0)|^2\). Les fonctions nulles au voisinage de zéro
sont denses dans \(L^2\) et y ont \(r\ge0\). Pourtant, pour une bosse
\(\phi(0)=1\), \(f_\eta(x)=\phi(x/\eta)\) donne
\(r(f_\eta)=\eta\|\phi\|_2^2-1<0\) pour \(\eta\) petit. L'évaluation
ponctuelle n'est pas continue en norme \(L^2\) ; elle l'est en \(H^1\)
en dimension un.

\textbf{Norme qui s'échappe.} \(r(f)=-\|f\|^2\),
\(R_n(f)=e^{2n}\|f\|^2\) et \(\varepsilon_n=e^{-n}\) satisfont
\(r(f)\ge-\varepsilon_nR_n(f)\) pour \(n\ge0\), avec
\(\varepsilon_n\to0\). Pourtant le signe reste négatif. Ce qui doit
tendre vers zéro est le produit d'erreur pour chaque test fixé.

\hypertarget{protocole-actif-crituxe8res-de-duxe9cision-et-portuxe9e}{%
\section{7. Protocole actif, critères de décision et
portée}\label{protocole-actif-crituxe8res-de-duxe9cision-et-portuxe9e}}

\begin{longtable}[]{@{}
  >{\raggedright\arraybackslash}p{(\columnwidth - 4\tabcolsep) * \real{0.3333}}
  >{\raggedright\arraybackslash}p{(\columnwidth - 4\tabcolsep) * \real{0.3333}}
  >{\raggedright\arraybackslash}p{(\columnwidth - 4\tabcolsep) * \real{0.3333}}@{}}
\toprule\noalign{}
\begin{minipage}[b]{\linewidth}\raggedright
Étape
\end{minipage} & \begin{minipage}[b]{\linewidth}\raggedright
Produit exigé
\end{minipage} & \begin{minipage}[b]{\linewidth}\raggedright
État de cette passe
\end{minipage} \\
\midrule\noalign{}
\endhead
\bottomrule\noalign{}
\endlastfoot
Domaines et transfert & Continuité indépendante de l'identité, cœur et
exhaustion & Lemme A établi à partir de la borne publiée ; transfert
rappelé et prouvé \\
Normalisations & Égalité exacte entre forme publiée et forme
arithmétique du corpus sur le domaine choisi & À dériver avant
comparaison numérique \\
Couplages & Formule de \(\mathfrak D(f_i,f_j)\) à partir des vrais
coefficients, avec phases complexes & Obligation ouverte \\
Assemblage & Garde d'image et dette de Schur pour toute la grammaire &
Lemme C prouvé ; règle arithmétique ouverte \\
Calculs & Contre-modèles, polarisation complexe, pivots singuliers &
Contrôles exacts exécutés par le programme joint \\
Conclusion RH & Quantificateurs sur tous les rangs et supports, gardes
satisfaites & Non obtenue \\
\end{longtable}

Le programme joint ne calcule pas l'opérateur de Suzuki sur les premiers
: il vérifie exactement les identités algébriques et les pièges qui
décideront de la validité d'un tel calcul. Il ne fournit aucun
certificat numérique de RH.

\textbf{Décision de recherche.} Conserver l'identité de Gram comme cible
structurée et l'inégalité à erreur évanescente comme alternative. Le
prochain calcul utile doit exposer le défaut mixte avec ses conventions
et ses bornes d'erreur. Une simple extension de plafond ou une série de
normes positives n'apporterait pas ce renseignement.

Le Bellman documentaire conserve ses coûts et ses scores : aucune
probabilité de réussite n'est inventée. Ses investigations « couverture
constructive », « dette duale » et « test de paquets » reçoivent les
nouvelles pièces et les obligations précises. Les niveaux de progression
restent inchangés : un raccord fonctionnel ne ferme pas le signe global.

\hypertarget{ruxe9fuxe9rences-et-piuxe8ces-reliuxe9es}{%
\section{Références et pièces
reliées}\label{ruxe9fuxe9rences-et-piuxe8ces-reliuxe9es}}

\begin{itemize}
\tightlist
\item
  {[}SuzH23{]} M. Suzuki, \emph{On the Hilbert space derived from the
  Weil distribution}, arXiv:2301.00421, version consultée v3 ;
  propositions 1.2--1.3, théorème 1.4, §3.2.
  \href{https://arxiv.org/html/2301.00421v3}{Texte} ;
  \href{https://www.cambridge.org/core/journals/canadian-journal-of-mathematics/article/on-the-hilbert-space-derived-from-the-weil-distribution/FD33EBD117A6B448053450DD8D62ADB6}{publication}.
  Vérification ciblée le 15 octobre 2026.
\item
  Document 67, \emph{Construire des ponts avec Bellman}, ponts B2, B3,
  B5 : domaines, cœur et Schur.
\item
  Document 76, \emph{Weil : retrouver la piste du signe global},
  §2.3--2.4 : grammaires couvrantes et rôle de l'arithmétique.
\item
  Documents 78--79 : transport cohérent, compensations et obstructions
  aux résumés trop grossiers.
\item
  Programme :
  \texttt{3\_Preuves\_et\_calculs/coup\_double\_20261015/verifier\_coup\_double.py}
  ; résultats : \texttt{resultats.json} ; confrontation :
  \texttt{CONFRONTATION.md}.
\end{itemize}

\end{document}
